\documentclass[11pt,a4paper]{article}
\usepackage[margin=2.2cm]{geometry}
\usepackage[T1]{fontenc}
\usepackage{lmodern}
\usepackage{amsmath,amssymb}
\usepackage{graphicx}
\usepackage{booktabs}
\usepackage{array}
\usepackage{longtable}
\usepackage{caption}
\usepackage{float}
\usepackage[dvipsnames]{xcolor}
\usepackage[most]{tcolorbox}
\usepackage{url}
\def\UrlBreaks{\do\/\do-\do_\do.}
\usepackage[colorlinks=true,linkcolor=blue!50!black,urlcolor=blue!50!black,citecolor=blue!50!black]{hyperref}
\DeclareGraphicsExtensions{.pdf,.png}
\graphicspath{{figures/}}
\sloppy
\emergencystretch=2em
\captionsetup{font=small,labelfont=bf}

% ---------------------------------------------------------------- macros
\newcommand{\fpath}[1]{\nolinkurl{#1}}
\newcommand{\fcmd}[1]{\texttt{\detokenize{#1}}}
\newcommand{\pt}{\ensuremath{p_T}}
\newtcolorbox{choicebox}{colback=blue!5,colframe=blue!60!black,fonttitle=\bfseries,title=Choice,left=2pt,right=2pt,top=1pt,bottom=1pt}
\newtcolorbox{caveatbox}{colback=red!5,colframe=red!60!black,fonttitle=\bfseries,title=Caveat,left=2pt,right=2pt,top=1pt,bottom=1pt}
\newtcolorbox{decisionbox}{colback=orange!7,colframe=orange!80!black,fonttitle=\bfseries,title=Decision needed,left=2pt,right=2pt,top=1pt,bottom=1pt}
\newcommand{\choice}[1]{\begin{choicebox}\small #1\end{choicebox}}
\newcommand{\caveat}[1]{\begin{caveatbox}\small #1\end{caveatbox}}
\newcommand{\decision}[1]{\begin{decisionbox}\small #1\end{decisionbox}}

% script-generated numbers (scripts/make_qa.C): fallbacks, then values
\input{tables/qa_defaults.tex}
\input{tables/qa_numbers.tex}
\input{tables/dca_defaults.tex}
\input{tables/dca_numbers.tex}
\input{tables/effdca_defaults.tex}
\input{tables/effdca_numbers.tex}
\input{tables/ptres_defaults.tex}
\input{tables/ptres_numbers.tex}

\title{A simple event/track QA ntuple for sPHENIX Run-3 p+p\\
\texttt{DST\_TRKR\_TRACKS} (ana581\_2026p005\_v001)\\[4pt]
\large Tool description, first test and a DCA study}
\author{X.~Dong (dongxs) \quad -- \quad draft prepared with Hermes}
\date{\today}

\begin{document}
\maketitle

\begin{abstract}
\noindent
This note describes \fcmd{TrackQATree}, a small Fun4All module that reduces the Run-3 p+p track
DSTs to a flat ROOT ntuple with one tree of event variables (trigger/GL1 information, track and
vertex counts, crossing-0 vertex) and one tree of track variables (kinematics, fit quality,
cluster counts per detector, beam crossing, impact parameters to an associated vertex). As an
option the same job writes a trimmed DST that keeps only the track and vertex maps. On one DST
segment (run \runNo, segment 0, \nEvents{} events, \nTracks{} tracks) the ntuple takes
\sizeNtkB\,kB per event, $\ratioNt\times$ smaller than the input DST (\sizeDSTkB\,kB/event); the
trimmed DST takes \sizeTrimkB\,kB/event. The job needs \cpuMsPerEv\,ms of CPU per event, mostly
for reading the DST. All consistency checks pass, including an event-by-event comparison of the
track and vertex counts with the independent \fcmd{EventInfoTree} of the KFParticle test and an
entry-by-entry identity of the ntuple made from the trimmed DST. First QA distributions are shown.
A 100-segment sample (\sNEvents{} events, \sNTracks{} tracks) is used to study DCA$_{xy}$ and DCA$_z$
of tracks with $p>0.8$\,GeV/$c$ vs.\ $p$, $\eta$ and $\phi$: the core widths of tracks used in their
vertex fit are \sSigXyA{} and \sSigZA\,$\mu$m and nearly independent of $p$, the core means vary with
$\phi$ between \sMeanXyPhiMin{} and \sMeanXyPhiMax\,$\mu$m ($xy$) in the same way for both charges,
and the stored DCA uncertainties are \sPullXyA{} ($xy$) and \sPullZA{} ($z$) times too small.
The silicon-matching fraction of TPC tracks inside $|z|<10$\,cm, $|\eta|<1$ is \eEffAll\% and flat to
2.5\,GeV/$c$; it has \eLowPhiN{} azimuthal holes and drops to \eEffPtHigh\% at 6--10\,GeV/$c$, mostly from poorly
fitted TPC tracks and partly from the worse TPC-only momentum resolution (with quality cuts \rEffQualFour\% at
4--6\,GeV/$c$); the DCA does not depend on
the beam crossing, and the mean DCA$_{xy}$ rises with vertex $z$ by \eSlopeXyVz\,$\mu$m/cm.
\textbf{Limitations:} one run (1 and 100 segments); no event, trigger or track selection; the
impact-parameter widths are residuals to vertices that include the track and are not a
resolution; no dE/dx (needs geometry and clusters).
Boxes: \textcolor{blue!60!black}{blue} = choice made and why,
\textcolor{red!60!black}{red} = caveat / open point,
\textcolor{orange!80!black}{orange} = needs a decision by the analyser.
\end{abstract}

\tableofcontents
\clearpage

\section{Overview}
\label{sec:overview}

\paragraph{Goal.} A light-weight, fast, run-everywhere QA ntuple of the Run-3 p+p track DSTs: per
event the trigger (GL1) word, track and vertex counts; per track everything needed for basic
tracking QA (acceptance, cluster counts, fit quality, crossing, impact parameter). It complements
the KFParticle K$^0_S$/$\Lambda$/D$^0$ ntuples of the same project (\fpath{run3pp/analysis_note}),
whose pre-requisites (release, crossing structure, track types) are re-used here.

\paragraph{Findings on the test segment} (run \runNo, segment 0, \nEvents{} events):
\begin{itemize}\itemsep0pt
\item Output size: ntuple \sizeNtkB\,kB/event ($\ratioNt\times$ smaller than the DST,
  \bytesPerTrk\,B per track), trimmed DST \sizeTrimkB\,kB/event ($\ratioTrim\times$ smaller).
  Clusters and seeds are \pctClus\% and \pctSeeds\% of the DST.
\item Speed: \cpuMsPerEv\,ms CPU per event for the whole job, of which the module itself takes
  \modMsPerEv\,ms; a 1k-event segment runs in \wallTime{} (m:ss) wall time.
\item Content: \trkPerEv{} tracks per event, \fracINTT\% with INTT clusters, \fracMVTXonly\% MVTX
  without INTT, \fracTpcOnly\% TPC-only; \vtxPerEv{} vertices per event, of which \vtxXzeroPerEv{}
  at crossing 0; \fracEvVtxXzero\% of the events have a crossing-0 vertex.
\item Validation: \nFailed{} of the \nChecks{} consistency checks of Sec.~\ref{sec:audit} fail, including
  event-by-event equality of six track/vertex counters with the KFParticle-test \fcmd{EventInfo}
  tree (\nKFPAgree/\nEvents{} events) and an identical ntuple when re-made from the trimmed DST.
\item Vertex association: a track's stored vertex id always points to a vertex of its own
  crossing; TPC-only tracks never carry a vertex id. \fracMatchOne\% of the tracks were used in
  their vertex fit, \fracMatchTwo\% carry a vertex id without being used in the fit,
  \fracMatchThree\% (all TPC-only) are associated with the closest crossing-0 vertex in $z$, and
  \fracMatchZero\% have no vertex in their crossing.
\end{itemize}

\paragraph{Silicon efficiency and DCA dependences} (Sec.~\ref{sec:effdca}, same 100 segments):
\begin{itemize}\itemsep0pt
\item silicon-matched fraction of TPC tracks (INTT\,$\geq1$, MVTX\,$\geq2$) in the matching acceptance
  $|z_{\rm PCA}|<10$\,cm, $|\eta|<1.5$ ($|\eta|<1$ vs.\ $p_T$): \eEffAll\% ($|\eta|<1$), \eEffHi\% ($p_T>1$),
  flat to 2.5\,GeV/$c$, \eEffPtHigh\% at 6--10\,GeV/$c$ (Sec.~\ref{sec:ptres}: mostly poorly fitted TPC tracks plus
  resolution migration; \rEffQualFour\% at 4--6\,GeV/$c$ with quality cuts); \eLowPhiN{} $\phi$ bins below 70\%;
  flat in time, bunch number and crossing-0 multiplicity;
\item DCA widths independent of beam crossing and time; mean DCA$_{xy}$ vs.\ vertex $z$ slope \eSlopeXyVz\,$\mu$m/cm.
\end{itemize}

\paragraph{DCA study} (Sec.~\ref{sec:dca}, run 79509 segments 0--99, \sNEvents{} events, ntuple
saved as \fpath{ntuples/TrackQA_run79509_seg000-099.root}), tracks with MVTX\,$\geq3$, INTT\,$\geq1$,
TPC\,$\geq20$, $p>0.8$\,GeV/$c$, used in their vertex fit:
\begin{itemize}\itemsep0pt
\item core $\sigma_{xy}=\sSigXyA\,\mu$m, $\sigma_z=\sSigZA\,\mu$m; $\sigma_{xy}$ \sSigXyAlowP$\to$\sSigXyAhighP\,$\mu$m
  from 0.8 to 10\,GeV/$c$ (no visible multiple-scattering term), $\sigma_z$ \sSigZEtaMin--\sSigZEtaMax\,$\mu$m over $|\eta|<1.2$;
\item $\phi$-dependent core means (DCA$_{xy}$ \sMeanXyPhiMin{} to \sMeanXyPhiMax\,$\mu$m, DCA$_z$ \sMeanZPhiMin{} to
  \sMeanZPhiMax\,$\mu$m), identical for both charges: likely residual silicon misalignment;
\item DCA pulls have core widths \sPullXyA{} ($xy$) and \sPullZA{} ($z$): stored uncertainties too small;
\item tracks carrying a vertex id without being in the fit are not prompt (sample B).
\end{itemize}

\paragraph{Choices} (each in a blue box where it is used):
no event or track selection in the ntuple (Sec.~\ref{sec:tool});
no geometry, CDB or reconstruction in the job (Sec.~\ref{sec:tool});
vertex for the DCA = stored vertex id, else closest same-crossing vertex (Sec.~\ref{sec:vtxassoc});
trimmed DST keeps \fcmd{SvtxTrackMap}, \fcmd{SvtxVertexMap}, GL1, event header and sync nodes
(Sec.~\ref{sec:size}); ZSTD-5 compression, 32-bit floats and 8/16-bit integers for counters
(Sec.~\ref{sec:ntuple}).

\section{Dataset, software and computing}
\label{sec:dataset}

\begin{table}[H]
\centering\small
\begin{tabular}{>{\raggedright}p{0.25\textwidth}p{0.68\textwidth}}
\toprule
input & \fpath{/sphenix/lustre01/sphnxpro/production/run3pp/physics/ana581_2026p005_v001/DST_TRKR_TRACKS/run_00079500_00079600/DST_TRKR_TRACKS_run3pp_ana581_2026p005_v001-00079509-00000.root}\\
events in input (tree \fcmd{T}) & \nDSTEntries \\
input size & \sizeDSTMB\,MB (\sizeDSTkB\,kB/event; 1\,kB = 1024\,B throughout) \\
release & \fcmd{ana.581} (\fpath{/cvmfs/sphenix.sdcc.bnl.gov/alma9.2-gcc-14.2.0/release/release_ana/ana.581}), ROOT 6.32.06, gcc 14.2.0, via \fpath{run3pp/setup_env.sh} \\
module library & \fpath{/gpfs02/sphenix/user/dongxs/run3pp/install/lib/libtrackqatree.so} \\
job & interactive on \fcmd{sphnxuser01}, \fpath{trackqa/macro/run_trackqa.sh 0 <DST> <outdir> 1} \\
CPU / wall / memory & user \cpuUser\,s + system \cpuSys\,s, wall \wallTime, max RSS \rssMB\,MB \\
module time & \modMsPerEv\,ms/event (Fun4All timer of \fcmd{TrackQATree}) \\
outputs & \fpath{/gpfs02/sphenix/user/dongxs/run3pp/trackqa/output/test79509/}; local copies in \fpath{trackqa/ntuples/} (md5 checked) \\
\bottomrule
\end{tabular}
\caption{Test sample and job. Source: \fpath{logs/dataset_software.txt} (remote \fcmd{ls}, release
banner, DST branch sizes) and \fpath{logs/run_trackqa_79509_00000.log} (job log with
\fcmd{/usr/bin/time -v}). The same segment was the first file of the KFParticle test sample.}
\label{tab:dataset}
\end{table}

The run was chosen because it is the KFParticle test run, so the event counters can be compared
event by event with that independent code (Sec.~\ref{sec:audit}). Run 79509 is a light run; busier
runs (e.g.\ 80906) have larger DSTs and the time per event of this read-dominated job will scale
with the DST size (not measured here).

\section{The tool}
\label{sec:tool}

\paragraph{Files} (local mirror \fpath{~/Work/Hermes/sPHENIX/run3pp/trackqa/}, cluster
\fpath{/gpfs02/sphenix/user/dongxs/run3pp/trackqa/}):
\begin{itemize}\itemsep0pt
\item \fpath{src/TrackQATree/} -- the SubsysReco module (autotools package, libraries
  \fcmd{fun4all}, \fcmd{ffaobjects}, \fcmd{ffarawobjects}, \fcmd{trackbase_historic(_io)},
  \fcmd{globalvertex_io}); built by \fpath{build.sh} into the project install area shared with
  \fcmd{EventInfoTree}.
\item \fpath{macro/Fun4All_TrackQA.C(nEvents, input, outdir, trimDST, minPt)} -- DST input
  manager (single file or a text list of files of one run), the module, and optionally a
  \fcmd{Fun4AllDstOutputManager} for the trimmed DST. Prints \fcmd{All done} only when all
  outputs exist.
\item \fpath{macro/run_trackqa.sh <nEvents> <input> <outdir> [trimDST 0/1] [minPt]} -- job wrapper:
  clean environment, scratch directory, log with \fcmd{/usr/bin/time -v}, fails unless the log
  contains \fcmd{All done}. Same conventions as \fcmd{run_kfp.sh}, so the existing condor
  scripts can be adapted by changing the executable.
\item \fpath{tools/check_vertex_id.C} -- diagnostic used to define the vertex association
  (Sec.~\ref{sec:vtxassoc}).
\item \fpath{analysis_note/scripts/make_qa.C} -- figures, numbers and checks of this note (runs
  locally on the downloaded ntuple).
\end{itemize}

\choice{The job reads the DST nodes only: no CDB, geometry, field map or track refit. The DCA
calculation (\fcmd{TrackAnalysisUtils::get_dca}) needs only the stored track state and
covariance at the point of closest approach. This keeps the job at \cpuMsPerEv\,ms/event and
independent of calibration tags. Quantities that need geometry or clusters (dE/dx, residuals)
are therefore not in the ntuple.}

\choice{No event or track selection is applied when writing: every DST event has an
\fcmd{event} entry and every track of \fcmd{SvtxTrackMap} a \fcmd{track} entry
(\fcmd{minPt}=0 by default). Selections are cheap to apply offline, and the
counts can be checked against the DST (Sec.~\ref{sec:audit}). The \fcmd{minPt} argument exists
to shrink the track tree if needed; tracks below it are still counted in the event tree.}

\subsection{Track classes and crossing}
Following the KFParticle study, tracks are classified by their cluster content
(\fcmd{SvtxTrack_v5} cluster counters): \emph{INTT} ($\geq1$ INTT cluster, crossing measured by
INTT timing), \emph{MVTX-only} (MVTX but no INTT clusters) and \emph{TPC-only} (no silicon
cluster). The track beam crossing \fcmd{crossing} is relative to the GL1 BCO of the event; TPC-only
tracks always carry crossing 0 because no silicon timing is available (verified for every track,
check \fcmd{chkTpcOnlyXzero}).

\subsection{Vertex association for the DCA}
\label{sec:vtxassoc}
\fcmd{SvtxTrack::get_vertex_id()} was checked on 200 events with \fpath{tools/check_vertex_id.C}:
when set, it always points to an existing vertex of the track's own crossing; the vertex track
lists are consistent with the ids (every listed track points back to the vertex), but many tracks
carry the id of a vertex they were not used in. The ntuple therefore stores the class
\fcmd{vtxMatch}:
\begin{center}\small
\begin{tabular}{clr}
\toprule
\fcmd{vtxMatch} & definition & tracks (\%)\\
\midrule
1 & vertex = stored vertex id, track in that vertex's track list (used in the fit) & \nMatchOne{} (\fracMatchOne)\\
2 & vertex = stored vertex id, track not in the fit & \nMatchTwo{} (\fracMatchTwo)\\
3 & no vertex id; vertex of the same crossing closest in $z$ to the track PCA & \nMatchThree{} (\fracMatchThree)\\
0 & no vertex in the track's crossing; DCA and vertex branches $=-999$ & \nMatchZero{} (\fracMatchZero)\\
\bottomrule
\end{tabular}
\end{center}
In the test segment class 3 contains only TPC-only tracks and classes 1--2 only silicon tracks
(check \fcmd{chkTpcVtxId}).

\choice{DCA reference = the stored vertex id when present (the tracking's own assignment),
otherwise the closest vertex in $z$ of the same crossing. The class is stored, so analyses can
choose; e.g.\ \fcmd{vtxMatch==1} for primary-track QA.}

\caveat{For \fcmd{vtxMatch==1} the track was used in the vertex fit, so its DCA is biased towards
zero, strongly for vertices with few tracks (Sec.~\ref{sec:qa}). For TPC-only tracks (class 3) the
vertex is a crossing-0 vertex by construction, while most TPC-only tracks come from other
crossings (their $z$ is then shifted by the drift-time offset): their DCA is not a primary/secondary
discriminator.}

\section{Ntuple content}
\label{sec:ntuple}

Output file \fpath{TrackQA_<run>_<seg>.root} with trees \fcmd{event} and \fcmd{track}
(ZSTD level 5). A track is linked to its event through \fcmd{ievt} (entry number in
\fcmd{event}) or through (\fcmd{run}, \fcmd{event}); e.g.\
\fcmd{track->AddFriend} is not possible (different entry counts), use
\fcmd{event->GetEntry(ievt)} or an index on (\fcmd{run},\fcmd{event}).
Ntuple version~2 (used from Sec.~\ref{sec:ptres} on, 100-segment sample) adds the TPC-seed
\fcmd{tpcSeedPt}, \fcmd{tpcSeedEta}, \fcmd{tpcSeedPhi} ($-999$ when the DST has no \fcmd{TpcTrackSeedContainer},
e.g.\ the trimmed DST). Table~\ref{tab:branches} (version~1, test segment) lists every other branch with its ROOT leaf type and its share of the
compressed size of its tree; the event tree is \evtTreekBPerEv\,kB/event and the track tree
\trkTreekBPerEv\,kB/event.

\choice{32-bit floats for all continuous variables, 8-bit unsigned integers for cluster counts
and 16-bit integers for charge, ndf, crossing and class flags. Momentum components are not
stored: $p_x = p_T\cos\phi$, $p_y = p_T\sin\phi$, $p_z = p_T\sinh\eta$.}

{\small
\begin{longtable}{>{\raggedright}p{0.17\textwidth}c>{\raggedright}p{0.58\textwidth}r}
\caption{Branches of the ntuple. Type: ROOT leaf code (I int32, L int64, l uint64, F float,
S int16, b uint8, O bool). Last column: share of the compressed size of the tree in \% (test
segment). Source: \fpath{tables/tab_branches.tex} written by \fpath{scripts/make_qa.C}.}
\label{tab:branches}\\
\toprule
branch & type & content & size (\%)\\
\endfirsthead
\toprule
branch & type & content & size (\%)\\
\midrule
\endhead
\bottomrule
\endfoot
\csname @@input\endcsname tables/tab_branches.tex

\end{longtable}
}

\section{Trimming: DST vs.\ trimmed DST vs.\ ntuple}
\label{sec:size}

\begin{table}[H]
\centering\small
\input{tables/tab_dstsize.tex}
\caption{Compressed size per event of each node of the production DST and of the trimmed DST
(test segment, \nEvents{} events; ROOT \fcmd{GetZipBytes} of the branches of tree \fcmd{T}, run tree
not included). The trimmed DST is written by \fcmd{Fun4AllDstOutputManager} with its default
compression, which makes \fcmd{SvtxTrackMap} slightly larger than in the input. Source:
\fpath{logs/dataset_software.txt}.}
\label{tab:dstsize}
\end{table}

Clusters (\fcmd{TRKR_CLUSTER_SEED}) are \pctClus\% and the three seed containers \pctSeeds\% of
the DST; the track map is \pctTrkMap\% (\trkMapkBPerEv\,kB/event) and the vertex map
\pctVtxMap\%. File sizes per event:

\begin{center}\small
\begin{tabular}{lrrr}
\toprule
file & size (MB) & kB/event & reduction\\
\midrule
production DST & \sizeDSTMB & \sizeDSTkB & 1\\
trimmed DST (track + vertex maps, GL1, header) & \sizeTrimMB & \sizeTrimkB & \ratioTrim\\
TrackQA ntuple & \sizeNtMB & \sizeNtkB & \ratioNt\\
\bottomrule
\end{tabular}
\end{center}

\choice{Trimmed DST = \fcmd{Sync}, \fcmd{EventHeader}, \fcmd{GL1RAWHIT}, \fcmd{SvtxTrackMap},
\fcmd{SvtxVertexMap}. It keeps the full \fcmd{SvtxTrack} objects (states, covariances, seed ids),
so modules that need only tracks and vertices can run on it; for \fcmd{TrackQATree} this was
verified: re-running \fcmd{TrackQATree} on it reproduces the
ntuple exactly (check \fcmd{chkTrim}). The file name keeps the production pattern
\fpath{DST_TRKR_TRACKS_TRIM-<run>-<seg>.root} so that \fcmd{Fun4AllUtils::GetRunSegment} works on it.}

\caveat{Anything that needs clusters or seeds (dE/dx, residuals, refits, KFParticle detailed
tracking or dE/dx branches) does not work on the trimmed DST. Running KFParticle on it was not
tested.}

\section{First QA distributions}
\label{sec:qa}

All figures: test segment (run \runNo, segment 0, \nEvents{} events), no event selection, made by
\fpath{scripts/make_qa.C} from \fpath{ntuples/TrackQA_00079509_00000.root}. Error bars, where shown,
are statistical. Statements marked \emph{(by eye)} are read from the figure and not computed.

\subsection{Event level}

\begin{figure}[H]
\centering
\includegraphics[width=0.6\textwidth]{fig_evt_ntrk}
\caption{Tracks per event: all tracks (black), tracks with at least one silicon cluster (blue),
TPC-only tracks (red) and silicon tracks at crossing 0 (green). Means: \trkPerEv{} all,
\silPerEv{} silicon, \tpcOnlyPerEv{} TPC-only, \silXzeroPerEv{} silicon at crossing 0 per event.
Most tracks come from other bunch crossings within the streaming readout window: these are frame totals,
not the multiplicity of one collision (for that see the same-crossing multiplicity in Sec.~\ref{sec:effdca}).}
\label{fig:ntrk}
\end{figure}

\begin{figure}[H]
\centering
\includegraphics[width=0.9\textwidth]{fig_evt_vtx}
\caption{Left: vertices per event, all (black) and at crossing 0 (red); means \vtxPerEv{} and
\vtxXzeroPerEv. Right: $z$ of the crossing-0 vertex with most tracks, for the
\fracEvVtxXzero\% of events that have one; mean \vzMean\,cm, RMS \vzRMS\,cm. Its mean transverse
position is $(\vxMean, \vyMean)\,\mu$m and it has on average \vntrkMean{} tracks.}
\label{fig:vtx}
\end{figure}

\begin{figure}[H]
\centering
\includegraphics[width=0.9\textwidth]{fig_evt_crossing}
\caption{Left: beam crossing of silicon tracks (relative to the GL1 BCO), per event. The spike at 0
is the triggered crossing; the plateau covers the streaming window (by eye: about $-100$ to $+480$).
Right: number of distinct crossings with at least one silicon track per event, mean \xingPerEv.}
\label{fig:xing}
\end{figure}

\subsection{Track level}

\begin{figure}[H]
\centering
\includegraphics[width=\textwidth]{fig_trk_kin}
\caption{Track $p_T$, $\eta$ and $\phi$ per event for INTT (blue), MVTX-only (green) and
TPC-only (red) tracks; $\eta$ and $\phi$ for $p_T>0.2$\,GeV/$c$ (\fracPtLow\% of all tracks
are below). By eye: the INTT-track $\phi$ distribution has dips near $\phi\approx0.2$ and
$\approx1.3$\,rad, and the MVTX-only tracks have an excess at $\phi\approx0.2$\,rad, consistent
with an INTT inefficiency there being partly recovered by MVTX-only tracks.}
\label{fig:kin}
\end{figure}

\begin{figure}[H]
\centering
\includegraphics[width=0.9\textwidth]{fig_trk_etaphi}
\caption{$\eta$--$\phi$ occupancy (linear $z$) of silicon (left) and TPC-only (right) tracks with
$p_T>0.2$\,GeV/$c$, all events. By eye: a low-occupancy band in silicon tracks at
$\phi\approx1.3$--1.4\,rad over the full $\eta$ range.}
\label{fig:etaphi}
\end{figure}

\begin{figure}[H]
\centering
\includegraphics[width=0.9\textwidth]{fig_trk_hits}
\caption{Left: clusters per track in MVTX (green), INTT (blue), TPC (red) and TPOT (magenta), all
tracks. Ranges: TPC \ntpcMin--\ntpcMax{} (mean \ntpcMean), MVTX and INTT at most \nmvtxMax{} and
\ninttMax, TPOT at most \ntpotMax{} (\fracTpot\% of tracks have TPOT clusters). No track has fewer than \ntpcMin{} TPC clusters, presumably a requirement of the production
tracking (not checked in its configuration). Right: $\chi^2$/ndf of silicon (blue) and TPC-only (red)
tracks, each normalised to unit area; means in the shown range \qualSilMean{} and \qualTpcMean.}
\label{fig:hits}
\end{figure}

\begin{figure}[H]
\centering
\includegraphics[width=0.9\textwidth]{fig_trk_dca}
\caption{Transverse (left) and longitudinal (right) DCA for $p_T>0.5$\,GeV/$c$, split by vertex
association class (Sec.~\ref{sec:vtxassoc}): silicon tracks used in the vertex fit (black, class 1),
silicon tracks carrying a vertex id but not used in the fit (blue, class 2), TPC-only tracks with
the closest crossing-0 vertex (red, class 3). Each curve normalised to unit area. Class 1 shows a
narrow core; class 2 is broad (secondaries, tracks from other collisions in the same crossing,
poorly measured tracks); class 3 is flat within $\pm1$\,cm in DCA$_{xy}$ (by eye), as expected for
tracks mostly from other crossings.}
\label{fig:dca}
\end{figure}

\begin{figure}[H]
\centering
\includegraphics[width=0.6\textwidth]{fig_trk_dcares}
\caption{Gaussian core width (iterative fit in $\pm2\sigma$) of DCA$_{xy}$ (circles) and DCA$_z$
(squares) vs.\ $p_T$ for class-1 tracks with MVTX\,$\geq3$, INTT\,$\geq1$, TPC\,$\geq20$ clusters,
all vertices (black, full) and vertices with $\geq10$ tracks (blue, open). Values in
Table~\ref{tab:dcares}.}
\label{fig:dcares}
\end{figure}

\begin{table}[H]
\centering\small
\input{tables/tab_dcares.tex}
\caption{DCA core widths of Fig.~\ref{fig:dcares} with fit uncertainties; $N$ = tracks in the bin.}
\label{tab:dcares}
\end{table}

\caveat{The widths of Fig.~\ref{fig:dcares} are residuals to a vertex that contains the track, not
the DCA resolution: the track pulls the vertex towards itself. With vertices of $\geq10$ tracks
(less pull) the widths at low $p_T$ are \emph{larger}
(Table~\ref{tab:dcares}), consistent with a significant bias (the two samples also differ in
vertex population, so this is not a clean measurement of the bias). A resolution needs an unbiased
vertex (refit without the track, or the vertex of the other tracks); not done here.}

\section{Silicon efficiency and DCA vs.\ kinematics, crossing, vertex and multiplicity}
\label{sec:effdca}

Sample: the 100-segment ntuple (run 79509, segments 0--99, \eNEvents{} events; production and merging described in Sec.~\ref{sec:dca}).
Numbers, table and figures from \fpath{scripts/qa_effdca.C}; checks \fcmd{chkEDup} (\chkEDup: unique event
numbers), \fcmd{chkENoEv} (\chkENoEv: every track finds its event) and \fcmd{chkECount} (\chkECount:
efficiency denominators equal the counted tracks).

\paragraph{Definitions.}
\begin{itemize}\itemsep0pt
\item \textbf{Silicon efficiency} $=N(\text{TPC}\geq20 \wedge \text{INTT}\geq1 \wedge \text{MVTX}\geq2)\,/\,N(\text{TPC}\geq20)$,
  over tracks of \fcmd{SvtxTrackMap} (every track has $\geq21$ TPC clusters) inside the matching acceptance
  \textbf{$|z_{\rm PCA}|<10$\,cm and $|\eta|<1.5$}; for the $p_T$ dependence \textbf{$|\eta|<1.0$}. $z_{\rm PCA}$ is the
  track $z$ at its point of closest approach, the only $z$ available for TPC-only tracks (no vertex in their
  crossing); for in-time tracks it is close to the vertex $z$. Denominators: \eNDen{} tracks ($|\eta|<1$, all
  $p_T$) and \eNDenHi{} ($|\eta|<1.5$, $p_T>1$\,GeV/$c$). The plots vs.\ event $v_z$ and track $z$ apply only
  $|\eta|<1.5$. Binomial uncertainties (\fcmd{TEfficiency}, Clopper--Pearson).
\item \textbf{DCA sample}: tracks passing the numerator cuts that were used in their vertex fit
  (\fcmd{vtxMatch==1}), \eNDca{} tracks (\eNDcaHi{} with $p_T>1$\,GeV/$c$). Core mean and $\sigma$ as in
  Sec.~\ref{sec:dca} (iterative Gaussian in $\pm2\sigma$, 10-$\mu$m bins), fit errors only.
\item Multiplicity: since an event is a streaming frame with many crossings, the multiplicity used is the
  number of silicon tracks in the \emph{same crossing} as the track (DCA), or at crossing 0 (efficiency, see
  Fig.~\ref{fig:effevent}); the event total is not used.
\item Dependences on $\eta$, $\phi$, crossing, $z$, multiplicity and time: $p_T>1$\,GeV/$c$; the $p_T$ and $p$
  dependences use all tracks ($p_T>0.2$ and $p>0.3$\,GeV/$c$). The DCA sample has no extra $z$/$\eta$ cut.
\end{itemize}

\choice{Matching acceptance $|z_{\rm PCA}|<10$\,cm, $|\eta|<1.5$ ($|\eta|<1$ vs.\ $p_T$), chosen by the
analyser: it removes most TPC tracks outside the silicon acceptance, including TPC-only tracks from other
crossings whose $z$ is shifted by the drift-time offset (Sec.~\ref{sec:ptres}).}

\caveat{This "efficiency" is the fraction of TPC tracks with a silicon match, not a tracking efficiency:
even inside the acceptance the denominator contains secondaries, out-of-time tracks that happen to have
$|z|<10$\,cm, and fakes. Its shape is meaningful for QA (holes, trends), its absolute value is an
approximate lower limit of the matching efficiency.}

\caveat{The track beam crossing is only measured with silicon; TPC-only tracks carry crossing 0 by
default. The efficiency vs.\ track crossing is therefore undefined (the denominator at crossing $\neq0$
would contain only numerator tracks). The efficiency is shown vs.\ the GL1 bunch number of the event
instead, and the numerator crossing distribution is shown for reference.}

\subsection{Silicon efficiency}

\begin{figure}[H]
\centering
\includegraphics[width=\textwidth]{fig_eff_kin}
\caption{Silicon efficiency inside $|z_{\rm PCA}|<10$\,cm vs.\ $p_T$ ($|\eta|<1$), $\eta$ and $\phi$
($p_T>1$\,GeV/$c$, $|\eta|<1.5$). Inclusive: \eEffAll\% ($|\eta|<1$, all $p_T$), \eEffHi\% ($|\eta|<1.5$,
$p_T>1$\,GeV/$c$). It falls from \eEffPtLow\% at 0.2--0.3 to
\eEffPtHigh\% at 6--10\,GeV/$c$ (Table~\ref{tab:effpt}; explained in Sec.~\ref{sec:ptres}); vs.\ $\eta$ it stays within
\eEffEtaMin--\eEffEtaMax\%. Vs.\ $\phi$ (5-degree bins) it ranges from \eEffPhiMin{} to \eEffPhiMax\%;
\eLowPhiN{} bins are below 70\%, centred at $\phi=$ \eLowPhiList\,rad.}
\label{fig:effkin}
\end{figure}

\begin{figure}[H]
\centering
\includegraphics[width=0.95\textwidth]{fig_eff_2d}
\caption{Silicon efficiency maps, $|z_{\rm PCA}|<10$\,cm. Left: $\eta$--$\phi$, $p_T>1$\,GeV/$c$; the deepest hole
($\phi\approx0.1$--0.3\,rad) covers $\eta\lesssim0.7$ only, the shallower bands (e.g.\ $\phi\approx-1.4$ and
1.3\,rad) the full $\eta$ range (by eye): azimuthal (and for the deep hole also $z$) sectors of the silicon
detectors. The event-$v_z$ panel of Fig.~\ref{fig:effevent} has no track-$z$ cut and is therefore lower. Right: $\eta$--$p_T$; the drop above $\approx3$\,GeV/$c$ is present at all $\eta$ (by eye).}
\label{fig:eff2d}
\end{figure}

\begin{table}[H]
\centering\small
\input{tables/tab_eff_pt.tex}
\caption{Silicon efficiency vs.\ $p_T$, $|z_{\rm PCA}|<10$\,cm, $|\eta|<1$ (Clopper--Pearson 68\% intervals; shown as 0.0 when below 0.05\%).}
\label{tab:effpt}
\end{table}

\begin{figure}[H]
\centering
\includegraphics[width=\textwidth]{fig_eff_event}
\caption{Silicon efficiency ($p_T>1$\,GeV/$c$, $|\eta|<1.5$, $|z_{\rm PCA}|<10$\,cm except for the two $z$ panels) vs.\ GL1 bunch number (\eBunchMin--\eBunchMax; range
\eEffBunchMin--\eEffBunchMax\%, the extremes are the sparsely populated bunches near the abort gap), event
$v_z$ of the crossing-0 vertex (events that have one; \eEffVzMin--\eEffVzMax\%), track $z$ at the PCA,
the number of silicon tracks at crossing 0 of the event, $N_{\rm Si}(X{=}0)$ (\eEffNtrkMin--\eEffNtrkMax\%; this is the multiplicity of the triggered collision only -- the TPC-only tracks of the denominator have no measured crossing, so a per-track same-crossing multiplicity cannot be defined for them) and event entry in the
merged file (time order, $\approx$1 segment per 1000 entries; \eEffEntryMin--\eEffEntryMax\%: stable over
the run). Track $z$: plateau of \eEffPcazPlateau\% for $|z|<10$\,cm, falling steeply beyond $|z|\approx15$\,cm,
the edge of the MVTX acceptance. Bottom right: crossing of the numerator tracks; \eFracNumXzero\% are in the
bin containing crossing 0.}
\label{fig:effevent}
\end{figure}

\subsection{DCA vs.\ momentum, $\eta$ and $\phi$}
Each figure has four panels: DCA$_{xy}$ and DCA$_z$ vs.\ the variable (log colour scale, $\pm0.1$\,cm shown,
40-$\mu$m bins), then the core $\sigma$ and core mean of DCA$_{xy}$ (black circles) and DCA$_z$ (red squares).

\begin{figure}[H]
\centering
\includegraphics[width=\textwidth]{fig_dcav_p}
\caption{DCA vs.\ $p$ (all $p_T$). $\sigma_{xy}$ falls from \eSigXyPfirst{} (0.3--0.4\,GeV/$c$) to
\eSigXyPlast\,$\mu$m (6--10\,GeV/$c$) and flattens above $\approx1$\,GeV/$c$; $\sigma_z$ stays within
\eSigZPMin--\eSigZPMax\,$\mu$m. Means within \eMeanXyPMin{} to \eMeanXyPMax\,$\mu$m ($xy$) and
\eMeanZPMin{} to \eMeanZPMax\,$\mu$m ($z$).}
\label{fig:dcavp}
\end{figure}

\begin{figure}[H]
\centering
\includegraphics[width=\textwidth]{fig_dcav_eta}
\caption{DCA vs.\ $\eta$, $p_T>1$\,GeV/$c$. $\sigma_{xy}$ flat (\eSigXyEtaMin--\eSigXyEtaMax\,$\mu$m),
$\sigma_z$ from \eSigZEtaMin{} (central) to \eSigZEtaMax\,$\mu$m (edges). Means: $xy$ \eMeanXyEtaMin{} to
\eMeanXyEtaMax\,$\mu$m, $z$ \eMeanZEtaMin{} to \eMeanZEtaMax\,$\mu$m, with an $\eta$-odd pattern in $z$ (by eye).}
\label{fig:dcaveta}
\end{figure}

\begin{figure}[H]
\centering
\includegraphics[width=\textwidth]{fig_dcav_phi}
\caption{DCA vs.\ $\phi$ (5-degree bins), $p_T>1$\,GeV/$c$. Core means $xy$ \eMeanXyPhiMin{} to
\eMeanXyPhiMax\,$\mu$m and $z$ \eMeanZPhiMin{} to \eMeanZPhiMax\,$\mu$m; widths $\sigma_{xy}$
\eSigXyPhiMin--\eSigXyPhiMax, $\sigma_z$ \eSigZPhiMin--\eSigZPhiMax\,$\mu$m. The structure is
step-like (sector by sector, by eye) and is visible as shifts of the band in the 2D panels.}
\label{fig:dcavphi}
\end{figure}

\subsection{DCA vs.\ crossing, vertex $z$, same-crossing multiplicity and time}

\begin{figure}[H]
\centering
\includegraphics[width=\textwidth]{fig_dcav_crossing}
\caption{DCA vs.\ track beam crossing (bins of 20), $p_T>1$\,GeV/$c$. $\sigma_{xy}$ is the same in the bin
containing crossing 0 (\eSigXyXbinZero\,$\mu$m) and in all other bins together (\eSigXyXbinOther\,$\mu$m);
over the bins $\sigma_{xy}$ ranges over \eSigXyXingMin--\eSigXyXingMax\,$\mu$m and $\sigma_z$ over
\eSigZXingMin--\eSigZXingMax\,$\mu$m, the lowest values in the sparse edge bins of the streaming window.
Means within \eMeanXyXingMin{} to \eMeanXyXingMax{} ($xy$) and \eMeanZXingMin{} to \eMeanZXingMax\,$\mu$m ($z$):
the DCA does not depend on the crossing.}
\label{fig:dcavx}
\end{figure}

\begin{figure}[H]
\centering
\includegraphics[width=\textwidth]{fig_dcav_vz}
\caption{DCA vs.\ $z$ of the associated vertex, $p_T>1$\,GeV/$c$. The mean DCA$_{xy}$ rises linearly with
$z$: slope $\eSlopeXyVz\pm\eESlopeXyVz\,\mu$m/cm (fit $|z|<15$\,cm), from \eMeanXyVzMin{} to
\eMeanXyVzMax\,$\mu$m, which looks like a relative tilt of the silicon barrels or a $z$-dependent
misalignment (interpretation not verified). $\sigma_{xy}$ \eSigXyVzMin--\eSigXyVzMax\,$\mu$m and $\sigma_z$
\eSigZVzMin--\eSigZVzMax\,$\mu$m, larger at negative $z$.}
\label{fig:dcavvz}
\end{figure}

\begin{figure}[H]
\centering
\includegraphics[width=\textwidth]{fig_dcav_ntrk}
\caption{DCA vs.\ the number of silicon tracks (MVTX or INTT cluster, any $p_T$) in the \emph{same beam crossing}
as the track, $N_{\rm Si}$, for $p_T>1$\,GeV/$c$ (log $x$). An event (streaming frame) contains on average
\eXingPerEv{} crossings with silicon tracks, with \eSilPerXing{} silicon tracks per such crossing (median
\eSilPerXingMedian), so the event total is not the multiplicity of the collision the track comes from.
$\sigma_{xy}$ rises from \eSigXyNtrkMin\,$\mu$m ($N_{\rm Si}=2$) and saturates at about 70\,$\mu$m (by eye) from
$N_{\rm Si}\approx6$; $\sigma_z$ from \eSigZNtrkMin{} to about 105\,$\mu$m (range \eSigZNtrkMin--\eSigZNtrkMax).
The rise is the vertex self-bias of small vertices (Fig.~\ref{fig:dcantrk}), not a multiplicity-dependent
resolution. Means within \eMeanXyNtrkMin{} to \eMeanXyNtrkMax{} ($xy$) and \eMeanZNtrkMin{} to \eMeanZNtrkMax\,$\mu$m ($z$),
the extreme values in the sparse last bin.}
\label{fig:dcavnt}
\end{figure}

\begin{figure}[H]
\centering
\includegraphics[width=\textwidth]{fig_dcav_entry}
\caption{DCA vs.\ event entry / 1000 ($\approx$ segment, time order within the run), $p_T>1$\,GeV/$c$:
$\sigma_{xy}$ \eSigXyEntryMin--\eSigXyEntryMax, $\sigma_z$ \eSigZEntryMin--\eSigZEntryMax\,$\mu$m,
means within \eMeanXyEntryMin{} to \eMeanXyEntryMax{} and \eMeanZEntryMin{} to \eMeanZEntryMax\,$\mu$m: stable
over the 100 segments.}
\label{fig:dcavent}
\end{figure}

\subsection{Summary of this section}
\begin{itemize}\itemsep0pt
\item Inside the matching acceptance ($|z_{\rm PCA}|<10$\,cm, $|\eta|<1$) the silicon efficiency is
  \eEffAll\% and flat ($\approx$90--92\%) from 0.2 to 2.5\,GeV/$c$; it is stable in time, bunch number and
  crossing-0 multiplicity, has \eLowPhiN{} azimuthal bins below 70\% (at $\phi=$ \eLowPhiList\,rad), falls
  steeply outside the MVTX $z$ acceptance ($|z|\gtrsim15$\,cm), and drops above 3\,GeV/$c$ to \eEffPtHigh\% at
  6--10\,GeV/$c$. Sec.~\ref{sec:ptres}: the remaining drop is mostly poorly fitted TPC tracks, plus a small
  momentum-resolution migration.
\item DCA widths do not depend on the crossing or on time; vs.\ the same-crossing multiplicity they only show the vertex self-bias below $N_{\rm Si}\approx6$. The DCA means depend on $\phi$ (up to
  $\pm$50--70\,$\mu$m) and the DCA$_{xy}$ mean on vertex $z$ (\eSlopeXyVz\,$\mu$m/cm), both pointing to
  residual silicon misalignment in this production.
\end{itemize}

\section{Why does the silicon efficiency drop at high $p_T$?}
\label{sec:ptres}

Inside the matching acceptance ($|z_{\rm PCA}|<10$\,cm, $|\eta|<1$, \rNAcc{} tracks; Sec.~\ref{sec:effdca})
the silicon efficiency falls from $\approx$\rEffFullOneTwo\% at 1--2\,GeV/$c$ to \rEffFullFour\% at 4--6 and
\rEffFullSix\% at 6--10\,GeV/$c$. Without the acceptance cut the drop was much larger (to $\approx$45\% at
6--10\,GeV/$c$), mostly from TPC tracks outside the silicon $z$ acceptance or from other crossings (see the
composition numbers below); the acceptance is now part of the definition. Two explanations are tested here on the same
100 segments, with ntuple version~2 (adds the TPC-seed kinematics \fcmd{tpcSeedPt}, \fcmd{tpcSeedEta},
\fcmd{tpcSeedPhi} for every track; same events and tracks as version~1, all other numbers of the note
unchanged when re-run). Script: \fpath{scripts/ptres_study.C}; check \fcmd{chkRSeed} (\chkRSeed: every
track has a TPC seed).

\paragraph{(1) Momentum-resolution migration.} Numerator tracks are fitted with silicon + TPC, the
TPC-only tracks in the denominator with the TPC alone. A worse resolution smears the steeply falling
spectrum (local power law $p_T^{-\rSlope}$ for 1--10\,GeV/$c$) more, moving more tracks to high $p_T$, so
the ratio drops even for a flat true efficiency. \emph{Test}: compute the efficiency with the \emph{same}
TPC-only estimator for numerator and denominator, the pT of the TPC seed.

\paragraph{(2) Denominator composition.} TPC tracks that cannot have silicon clusters (poorly fitted or
fake tracks, residual out-of-time tracks) may make up a larger share at high $p_T$. \emph{Test}: add simple
quality cuts (TPC\,$\geq30$ clusters, track $\chi^2$/ndf\,$<3$) to numerator and denominator.

\begin{figure}[H]
\centering
\includegraphics[width=0.95\textwidth]{fig_ptres_eff}
\caption{Left: silicon efficiency vs.\ $p_T$ ($|z_{\rm PCA}|<10$\,cm, $|\eta|<1$) with the final-fit $p_T$
(black, as in Fig.~\ref{fig:effkin}), with the TPC-seed $p_T$ for numerator and denominator (red open), and with
the TPC-seed $p_T$ and the quality cuts TPC\,$\geq30$, $\chi^2$/ndf\,$<3$ (blue squares). Points shifted by $\pm3$\% in $p_T$ for visibility.
Right: $p_T$ spectra (same acceptance) of Si-matched (black/red) and TPC-only (blue/magenta) tracks with final-fit (solid) and
TPC-seed (dashed) $p_T$, all normalised to the number of Si-matched tracks.}
\label{fig:ptreseff}
\end{figure}

\begin{figure}[H]
\centering
\includegraphics[width=\textwidth]{fig_ptres_res}
\caption{Same acceptance. Left: TPC-seed vs.\ final-fit $p_T$ of Si-matched tracks (log scales). Middle and right: core
$\sigma$ and mean of $p_T/p_T^{\rm seed}-1$ (a curvature difference) vs.\ final-fit $p_T$, for Si-matched
tracks (seed vs.\ Si+TPC fit, black) and TPC-only tracks (seed vs.\ TPC-only fit, blue). If the Si+TPC fit is
much better than the seed, the black points are the TPC-seed resolution: \rResSiLow\% at 0.2--0.3 and
\rResSiOne\% at 1--1.2\,GeV/$c$, growing to $\approx$15--20\% at 3--4\,GeV/$c$ (Table~\ref{tab:ptres}). The TPC-only fit differs from its seed by only
\rResTpcOne\% at 1\,GeV/$c$, i.e.\ it has about the seed resolution. Above 4\,GeV/$c$ the Gaussian core fit is
not meaningful (few tracks, broad non-Gaussian distribution with a tail to final-fit $p_T\gg$ seed $p_T$).}
\label{fig:ptresres}
\end{figure}

\begin{table}[H]
\centering\small
\input{tables/tab_ptres.tex}
\caption{Per $p_T$ bin: core $\sigma$ and mean of $p_T/p_T^{\rm seed}-1$ for Si-matched and TPC-only tracks
(\%, bins above 4\,GeV/$c$ unreliable); migration factor = Si-matched tracks counted with seed $p_T$ / with
final-fit $p_T$; silicon efficiency with final-fit $p_T$, with seed $p_T$, and with seed $p_T$ and the quality
denominator. Statistical uncertainties of the final-fit efficiencies in Table~\ref{tab:effpt}.}
\label{tab:ptres}
\end{table}

\paragraph{Results} (all inside $|z_{\rm PCA}|<10$\,cm, $|\eta|<1$).
\begin{itemize}\itemsep0pt
\item \textbf{The resolution difference is real}: the TPC seed differs from the Si+TPC fit by a core
  $\sigma$ of \rResSiLow\% at low $p_T$ growing with $p_T$ (Table~\ref{tab:ptres}), so the TPC-only resolution
  is at least this if the Si+TPC fit is better (not verified without truth). The resulting migration of
  Si-matched tracks (migration column) is a few \% up to 3\,GeV/$c$ and about 13\% at 4--6\,GeV/$c$.
\item \textbf{It explains a small part of the drop}: with the seed $p_T$ for both, the efficiency is
  \rEffSeedFour\% (instead of \rEffFullFour\%) at 4--6 and \rEffSeedSix\% (instead of \rEffFullSix\%) at
  6--10\,GeV/$c$, against \rEffSeedOneTwo\% at 1--2\,GeV/$c$.
\item \textbf{Most of the remaining drop is poorly fitted TPC tracks}: with TPC\,$\geq30$ and
  $\chi^2$/ndf\,$<3$ the efficiency is \rEffQualOneTwo\% at 1--2, \rEffQualFour\% at 4--6 and \rEffQualSix\% at
  6--10\,GeV/$c$. Each cut alone (seed $p_T$; 1--2 / 4--6 / 6--10\,GeV/$c$): $\chi^2$/ndf\,$<3$
  \rEffCutChiOneTwo/\rEffCutChiFour/\rEffCutChiSix\%, TPC\,$\geq30$ \rEffCutNtpcOneTwo/\rEffCutNtpcFour/\rEffCutNtpcSix\%.
\item Composition \emph{before} the acceptance cut (final-fit $p_T$ 1--2 vs.\ 4--10\,GeV/$c$): TPC-only tracks
  have $|z_{\rm PCA}|>15$\,cm in \rCompZTpcOneTwo\% / \rCompZTpcFourTen\% of the cases (Si-matched:
  \rCompZSiOneTwo\% / \rCompZSiFourTen\%) -- the acceptance issue that the $|z|<10$\,cm cut removes. Their mean
  $\chi^2$/ndf rises from \rCompChiTpcOneTwo{} to \rCompChiTpcFourTen{} and their mean number of TPC clusters falls
  from \rCompNtpcTpcOneTwo{} to \rCompNtpcTpcFourTen. Si-matched tracks at 4--10\,GeV/$c$ also have a large mean
  $\chi^2$/ndf (\rCompChiSiFourTen, vs.\ \rCompChiSiOneTwo{} at 1--2).
\end{itemize}

\caveat{The seed-$p_T$ test assumes that the TPC seeds of Si-matched and TPC-only tracks have the same
resolution. The quality cuts select in-time, well-measured tracks in both numerator and denominator, so the
quality-selected efficiency is closer to a silicon matching efficiency for primary tracks but is itself a
biased sample (e.g.\ the $\chi^2$ cut removes poorly matched Si+TPC tracks). The large $\chi^2$/ndf of
high-$p_T$ Si-matched tracks and the tail to final-fit $p_T\gg$ seed $p_T$ (Fig.~\ref{fig:ptresres}, left)
suggest wrong silicon--TPC matches or misalignment at high $p_T$; not studied further.}

\section{DCA study with 100 segments}
\label{sec:dca}

\paragraph{Sample.} Run 79509, segments 0--99 (the first segment is the test segment of the previous
sections): \sNJobs{} interactive jobs on \fcmd{sphnxuser01} (6 in parallel; the condor queue held
$\sim$52k idle jobs), \sNJobsOK{} ended with \fcmd{All done}. Merged with \fcmd{hadd} into
\fpath{ntuples/TrackQA_run79509_seg000-099.root} (version~2, \sSizeMB\,MB, \sSizekB\,kB/event; cluster copy in
\fpath{trackqa/output/stat100_79509/merged/}, md5 checked): \sNEvents{} events and \sNTracks{}
tracks, $100\times$ the test segment. Job summary: \fpath{logs/stat100_jobs.txt}; checks
\fcmd{chkSJobs} (\chkSJobs) and \fcmd{chkSEntries} (\chkSEntries: merged entries = sum of the
per-job module counts). Numbers, tables and figures: \fpath{scripts/dca_study.C}.

\paragraph{Selection.} Tracks with MVTX\,$\geq3$, INTT\,$\geq1$, TPC\,$\geq20$ clusters and
$p=p_T\cosh\eta>0.8$\,GeV/$c$ (\sNGoodP{} tracks). They are split by vertex association
(Sec.~\ref{sec:vtxassoc}): \textbf{sample A} = used in the fit of their vertex (\fcmd{vtxMatch==1},
\sNSelA{} tracks), \textbf{sample B} = carrying the vertex id but not used in the fit
(\fcmd{vtxMatch==2}, \sNSelB). Only \sFracSelXzero\% of sample A are at crossing 0; the core widths
at crossing 0 (\sSigXyXzero/\sSigZXzero\,$\mu$m for $xy$/$z$) and at other crossings
(\sSigXyXother/\sSigZXother\,$\mu$m) agree, so all crossings are used.

\paragraph{Method.} Core width $\sigma$ and mean $\mu$ from a Gaussian fitted iteratively in
$\mu\pm2\sigma$ on 10-$\mu$m bins; uncertainties are fit errors (statistical). The DCA is
computed by \fcmd{TrackAnalysisUtils::get_dca} from the track state stored on the DST, with respect
to the vertex position from the DST.

\caveat{Sample A tracks are part of their vertex, so its DCA is a biased residual, not the
impact-parameter resolution. Fig.~\ref{fig:dcantrk} shows how strong the bias is for small
vertices; for $\geq6$ tracks the widths are almost constant, so the binned results are shown for all
vertices and for vertices with $\geq10$ tracks.}

\begin{figure}[H]
\centering
\includegraphics[width=0.9\textwidth]{fig_dca_incl}
\caption{DCA$_{xy}$ (left) and DCA$_z$ (right) of sample A (black) and B (blue), each normalised to
unit area, $\pm2$\,cm. Sample A: core $\sigma_{xy}=\sSigXyA\,\mu$m, $\sigma_z=\sSigZA\,\mu$m, means
$\sMeanXyA\pm\sEMeanXyA$ and $\sMeanZA\pm\sEMeanZA\,\mu$m; \sTailXyA\% ($xy$) and \sTailZA\% ($z$) of the
tracks lie beyond $5\sigma$ of the core (secondaries and badly measured tracks). Sample B has no
prompt core: \sFarFracXyB\% have $|{\rm DCA}_{xy}|>0.5$\,cm and \sFarFracZB\% have $|{\rm DCA}_z|>0.5$\,cm,
i.e.\ mostly tracks from other collisions in the same crossing, which carry the id of the wrong
vertex.}
\label{fig:dcaincl}
\end{figure}

\begin{figure}[H]
\centering
\includegraphics[width=0.9\textwidth]{fig_dca_core}
\caption{Core region of sample A with the Gaussian fit (red, solid in the fit range, dashed
extrapolation). Non-Gaussian tails start at about $\pm3\sigma$ (by eye).}
\label{fig:dcacore}
\end{figure}

\begin{figure}[H]
\centering
\includegraphics[width=0.9\textwidth]{fig_dca_vs_ntrk}
\caption{Sample A: core $\sigma$ (left) and mean (right) of DCA$_{xy}$ (circles) and DCA$_z$
(squares) vs.\ the number of tracks in the vertex. Two-track vertices: $\sigma_{xy}=\sSigXyNlow\,\mu$m,
$\sigma_z=\sSigZNlow\,\mu$m; highest bin: \sSigXyNhigh{} and \sSigZNhigh\,$\mu$m. The width first
increases (vertex bias falling) and then decreases slowly (better vertex).}
\label{fig:dcantrk}
\end{figure}

\subsection{Dependence on $p$, $\eta$ and $\phi$}

In Figs.~\ref{fig:dcap}--\ref{fig:dcaphi} black full markers are sample A with all vertices, blue
open markers sample A with vertices of $\geq10$ tracks; circles DCA$_{xy}$, squares DCA$_z$. All plots
have $p>0.8$\,GeV/$c$. Values for $p$ and $\eta$ in Tables~\ref{tab:dcap} and~\ref{tab:dcaeta}.

\begin{figure}[H]
\centering
\includegraphics[width=0.9\textwidth]{fig_dca_vs_p}
\caption{Core $\sigma$ (left) and mean (right) vs.\ $p$. $\sigma_{xy}$ goes from \sSigXyAlowP{} to
\sSigXyAhighP\,$\mu$m and $\sigma_z$ from \sSigZAlowP{} to \sSigZAhighP\,$\mu$m between the first
(0.8--1\,GeV/$c$) and last (6--10\,GeV/$c$) bins: no $1/p$ multiple-scattering term is visible above
0.8\,GeV/$c$, the width is dominated by a momentum-independent contribution.}
\label{fig:dcap}
\end{figure}

\begin{figure}[H]
\centering
\includegraphics[width=0.9\textwidth]{fig_dca_vs_eta}
\caption{Core $\sigma$ and mean vs.\ $\eta$. $\sigma_z$ (sample A, all vertices) ranges from
\sSigZEtaMin{} (central) to \sSigZEtaMax\,$\mu$m (edges), as expected from the longer extrapolation
and shallower crossing angle at large $|\eta|$; $\sigma_{xy}$ is flat within
\sSigXyEtaMin--\sSigXyEtaMax\,$\mu$m. The means vary within \sMeanZEtaMin{} to \sMeanZEtaMax\,$\mu$m
($z$); the DCA$_{xy}$ mean changes sign from positive at $\eta<-0.5$ to negative at $\eta>0$ (by eye).}
\label{fig:dcaeta}
\end{figure}

\begin{figure}[H]
\centering
\includegraphics[width=0.9\textwidth]{fig_dca_vs_phi}
\caption{Core $\sigma$ and mean vs.\ $\phi$ (36 bins). The mean DCA$_{xy}$ varies between
\sMeanXyPhiMin{} and \sMeanXyPhiMax\,$\mu$m (all vertices; \sMeanXyPhiMinTen{} to
\sMeanXyPhiMaxTen\,$\mu$m for $\geq10$ tracks), the mean DCA$_z$ between \sMeanZPhiMin{} and
\sMeanZPhiMax\,$\mu$m, with errors of about 1\,$\mu$m. $\sigma_{xy}$ ranges over
\sSigXyPhiMin--\sSigXyPhiMax\,$\mu$m and $\sigma_z$ over \sSigZPhiMin--\sSigZPhiMax\,$\mu$m, with
narrow peaks in $\sigma_z$ at several $\phi$ values (e.g.\ near $-2.0$, $-1.4$ and 1.3\,rad, by eye).
A first-harmonic fit $a_0+a_1\cos\phi+b_1\sin\phi$ to the DCA$_{xy}$ mean gives an amplitude of
\sModAmpA\,$\mu$m with $\chi^2/{\rm ndf}=\sModChiA$: the structure is not a simple vertex/beam
offset but has higher harmonics.}
\label{fig:dcaphi}
\end{figure}

\begin{figure}[H]
\centering
\includegraphics[width=0.55\textwidth]{fig_dca_phi_charge}
\caption{Core mean DCA$_{xy}$ vs.\ $\phi$ for negative (red, open) and positive (blue, full) tracks,
sample A, all vertices (points shifted by $\pm0.02$\,rad for visibility). The pattern is the same for
both charges: the largest difference in a bin is \sChargeMaxDiff\,$\mu$m
($\chi^2/{\rm ndf}=\sChargeChiNdf$ for the difference over \sChargeNbins{} bins, so significant but small).
A charge-independent pattern points to a geometric origin (silicon alignment), not to a
momentum/field effect.}
\label{fig:dcacharge}
\end{figure}

\begin{figure}[H]
\centering
\includegraphics[width=\textwidth]{fig_dca_2d}
\caption{Sample A, log colour scale: DCA$_{xy}$ vs.\ $\phi$, DCA$_z$ vs.\ $\phi$ and DCA$_z$ vs.\
$\eta$. The $\phi$ structure of Fig.~\ref{fig:dcaphi} is visible as steps of the band; DCA$_z$
widens towards large $|\eta|$.}
\label{fig:dca2d}
\end{figure}

\begin{figure}[H]
\centering
\includegraphics[width=0.9\textwidth]{fig_dca_pull}
\caption{DCA divided by its uncertainty from the track covariance, sample A (black) and B (blue),
with a unit Gaussian (red dashed). Sample A core widths: \sPullXyA{} ($xy$) and \sPullZA{} ($z$);
the stored uncertainties are smaller than the observed spread by these factors.}
\label{fig:dcapull}
\end{figure}

\begin{table}[H]
\centering\small
\input{tables/tab_dca_p.tex}
\caption{Core $\sigma$ and mean $\mu$ ($\mu$m, fit errors) vs.\ $p$; sample A. Errors below 0.5 show as 0.}
\label{tab:dcap}
\end{table}

\begin{table}[H]
\centering\small
\input{tables/tab_dca_eta.tex}
\caption{As Table~\ref{tab:dcap}, vs.\ $\eta$.}
\label{tab:dcaeta}
\end{table}

\subsection{Interpretation}
\begin{itemize}\itemsep0pt
\item The DCA core is $\phi$-dependent at the level of the core width itself: the mean offsets have
  an RMS over $\phi$ of \sMeanXyPhiRMS\,$\mu$m ($xy$) and \sMeanZPhiRMS\,$\mu$m ($z$), the same for both
  charges. Added in quadrature to per-$\phi$ widths of \sSigXyPhiMin--\sSigXyPhiMax\,$\mu$m this
  explains a good part of the inclusive \sSigXyA\,$\mu$m and its flatness vs.\ $p$. This looks like
  residual silicon misalignment in this production (not proven here).
\item The track DCA uncertainties underestimate the observed spread by factors \sPullXyA{} ($xy$) and \sPullZA{} ($z$); a DCA
  significance cut on these files would select far too many tracks as displaced.
\item Sample B is not a prompt sample; the stored vertex id of a track that was not used in the
  vertex fit should not be used for impact-parameter selections.
\end{itemize}

\caveat{All this is one run and biased residuals (sample A). The $\phi$ pattern is a statement about
residuals; a refit of the vertex without the track (needs geometry) would turn it into a resolution
and offset measurement. Statements marked \emph{(by eye)} are not computed.}

\section{Audit}
\label{sec:audit}

\subsection{Checks that passed}
All checks are computed by \fpath{scripts/make_qa.C} on the full test segment (\nEvents{} events,
\nTracks{} tracks) and written as PASS/FAIL into \fpath{tables/qa_numbers.tex} and
\fpath{logs/qa_numbers.txt}; \nFailed{} of \nChecks{} fail. The 100-segment checks are computed by \fpath{scripts/dca_study.C}; \sNFailed{} of \sNChecks{} fail.

\begin{center}\small
\begin{tabular}{>{\raggedright}p{0.17\textwidth}p{0.66\textwidth}c}
\toprule
check & what is compared & result\\
\midrule
\fcmd{chkEvEntries} & entries of tree \fcmd{event} = entries of the DST tree \fcmd{T} (\nDSTEntries) = events in the job log & \chkEvEntries\\
\fcmd{chkTrkEntries} & $\sum$ \fcmd{nTrk} over events = entries of tree \fcmd{track} = tracks in the job log (\fcmd{minPt}=0) & \chkTrkEntries\\
\fcmd{chkTypeSum} & INTT + MVTX-only + TPC-only = \fcmd{nTrk}, INTT + MVTX-only = \fcmd{nTrkSil}, in every event & \chkTypeSum\\
\fcmd{chkTpcOnlyXzero} & every TPC-only track has crossing 0 & \chkTpcOnlyXzero\\
\fcmd{chkTpcVtxId} & no TPC-only track in classes 1--2, no silicon track in class 3 & \chkTpcVtxId\\
\fcmd{chkMatchLog} & \fcmd{vtxMatch} counts in the ntuple = counts printed by the module & \chkMatchLog\\
\fcmd{chkDcaFinite} & DCA finite for every track with an associated vertex & \chkDcaFinite\\
\fcmd{chkSiSeed} & every silicon track has a silicon seed (id) & \chkSiSeed\\
\fcmd{chkTrim} & ntuple re-made from the trimmed DST identical entry by entry (event, id, $p_T$, $\eta$, $\phi$, $\chi^2$, crossing, cluster counts, class, DCAs; event counters, BCO, $v_z$); \nTrimDiff{} differences & \chkTrim\\
\fcmd{chkSJobs} & 100-segment sample: \sNJobs{} jobs, all with \fcmd{All done} (\fpath{scripts/dca_study.C}, \fpath{logs/stat100_jobs.txt}) & \chkSJobs\\
\fcmd{chkSEntries} & merged 100-segment ntuple: event and track entries = sum of per-job module counts & \chkSEntries\\
\fcmd{chkEDup} & 100-segment sample: event numbers unique (\fpath{scripts/qa_effdca.C}) & \chkEDup\\
\fcmd{chkENoEv} & 100-segment sample: every track finds its event by event number & \chkENoEv\\
\fcmd{chkECount} & efficiency denominators (all $p_T$, $p_T>1$) = counted tracks & \chkECount\\
\fcmd{chkRSeed} & every track of the 100-segment v2 ntuple has a TPC seed (\fpath{scripts/ptres_study.C}) & \chkRSeed\\
\fcmd{chkKFP} & per event \fcmd{nTrk}, \fcmd{nTrkINTT}, \fcmd{nTrkMVTXonly}, \fcmd{nTrkTPConly}, \fcmd{nTrkX0}, \fcmd{nVtx} equal to the \fcmd{EventInfoTree} of the KFParticle test (independent module, other job, matched on run and event number): \nKFPAgree{} of \nEvents{} events agree, \nKFPFound{} found & \chkKFP\\
\bottomrule
\end{tabular}
\end{center}

Vertex-id diagnostic (\fpath{tools/check_vertex_id.C}, first 200 events, output
\fpath{logs/check_vertex_id_79509_00000.txt}): every set vertex id exists in the vertex map and has
the track's crossing; every track listed in a vertex points back to it.

\subsection{Not verified (assumptions)}
\begin{itemize}\itemsep0pt
\item Only run 79509 (a light run) was processed (1 and 100 segments); other runs, busier DSTs and
  multi-file list input were not run with this macro (the list input code is the same as in the
  KFParticle macro, where it was tested).
\item The TPC-only $p_T$ resolution itself (only its difference to the Si+TPC fit is measured), and the origin of the large $\chi^2$/ndf of high-$p_T$ Si-matched tracks.
\item The misalignment interpretation of the $\phi$-dependent DCA means (Sec.~\ref{sec:dca}) is not proven.
\item The condor wrapper usage was not tested (only interactive runs).
\item \fcmd{TrackAnalysisUtils::get_dca} was used as is (sign convention and uncertainties not
  checked against an independent calculation).
\item Statements marked \emph{(by eye)} in Sec.~\ref{sec:qa}.
\item The minimum of \ntpcMin{} TPC clusters is attributed to the production tracking without
  checking its configuration.
\end{itemize}

\section{Open issues and next steps}
\label{sec:open}

Ordered by importance; items for the analyser tagged [decision].
\begin{enumerate}\itemsep2pt
\item[{[decision]}] \textbf{Scope of a production.} At \sizeNtkB\,kB/event the full run3pp sample
  ($\sim5\times10^8$ events, number from the KFParticle note) would need \projNtFullTB\,TB (1\,TB = $1024^4$\,B) and the trimmed DST \projTrimFullTB\,TB; a 10\%
  sample (\texttt{seg \% 10 == 0}, as for KFParticle) \projNtTenPctTB\,TB of ntuples. Run 79509 is a
  light run, so these are likely lower estimates. Options: per-run QA on a few
  segments per run, the 10\% sample, or a \fcmd{minPt} cut.
\decision{How many segments per run to process, and whether to write the trimmed DST
(\sizeTrimkB\,kB/event, \ratioTrimNt\,$\times$ the ntuple) or only the ntuple.}
\item \textbf{Unbiased DCA resolution and $\phi$-dependent offsets}: refit the vertex without the track
  before quoting resolutions; map the $\phi$ (and $\eta$) dependent DCA means of Sec.~\ref{sec:dca} to
  MVTX staves / INTT ladders and report them to the alignment group; check other runs.
\item \textbf{DCA uncertainties} are \sPullXyA{}/\sPullZA{} ($xy$/$z$) times too small: do not use
  DCA-significance cuts on these DSTs without a scale factor.
\item \textbf{dE/dx and per-layer information}: need geometry and clusters (production DST, not
  the trimmed one); could be added as an option with \fcmd{TrackingInit} as in the KFParticle macro.
\item \textbf{Run-by-run trends}: add a small script that fills per-run summary histograms
  (tracks per event by class, crossing-0 vertex fraction, $\phi$ holes) from many ntuples.
\item Test the condor wrapper on a few segments of other runs, including a busy run.
\end{enumerate}


\end{document}
